\documentclass[11pt]{letter}
\usepackage[margin=1in]{geometry}
\usepackage{hyperref}

\begin{document}

\begin{letter}{
Editor-in-Chief \\
Applied Soft Computing \\
}

\opening{Dear Editor,}

We submit our manuscript \textbf{``One Model, Three Tasks: Multi-Task Fine-Tuning Provides Implicit Regularization Against Label Hallucination in SOC Classification''} for consideration in \emph{Applied Soft Computing}.

\textbf{Summary.} We compare multi-task (MT) and single-task (ST) fine-tuning for SOC alert classification and discover that multi-task training provides implicit regularization against label hallucination. By learning a structured output format across three tasks (Classification, Severity, Attack Category), the MT model constrains vocabulary selection and reduces hallucinated labels by 1.39$\times$ on average and up to 33$\times$ at specific seeds. MT outperforms ST by 5.7 percentage points on the hardest task (Attack Category, $H=1.24$) while requiring only one model and one inference pass.

\textbf{Key contributions:}
\begin{enumerate}
\item Discovery of format regularization: structured multi-task output acts as an anti-hallucination mechanism.
\item A formal definition showing MT produces fewer out-of-vocabulary labels than ST.
\item A 33$\times$ error reduction at Seed 999, demonstrating the practical magnitude of the effect.
\end{enumerate}

\textbf{Relevance to Applied Soft Computing.} Multi-task learning is a core soft computing technique. Our finding that multi-task training provides a previously unknown benefit---hallucination suppression via format constraints---extends the theoretical understanding of why multi-task learning works and when it should be preferred.

\textbf{Novelty.} Prior multi-task work focuses on accuracy gains through shared representations. We are the first to identify multi-task training as an anti-hallucination mechanism and to show that the structured output format itself acts as an implicit constraint on the model's vocabulary.

This manuscript has not been published previously and is not under consideration elsewhere.

\closing{Respectfully submitted,\\[2em]
Nutthakorn Chalaemwongwan\\
Department of Computer Engineering, KOSEN-KMITL\\
King Mongkut's Institute of Technology Ladkrabang\\
nutthakorn.ch@kmitl.ac.th
}

\end{letter}
\end{document}
